\section{Theorem-Backed Reductions and Assumption Structure}
\label{app:theorem-backed}

This appendix gives a paper-level description of the most general reduction
object formalized in Lean. The point is to separate \emph{backend choices} from
\emph{mathematical obligations}. Once this is done, LLM summaries, learned
latent sketches, and supplied codecs all fall under the same three-law
interface.

\subsection{General State-Space View}

Let $\mathcal{S}$ be a \emph{theorem domain}: the objects actually manipulated
by the reduction procedure. A theorem-backed reduction consists of
\begin{itemize}[nosep]
    \item a leaf map $e:\Strings\to\mathcal{S}$,
    \item a merge map $m:\mathcal{S}\times\mathcal{S}\to\mathcal{S}$,
    \item a re-summary map $r:\mathcal{S}\to\mathcal{S}$,
    \item a semantic readout $\psi:\mathcal{S}\to\Y$.
\end{itemize}

The exact OPS-style obligations on this state space are:
\begin{align*}
\text{(TB1)}\qquad &\psi(e(x)) = \fstar(x), \\
\text{(TB2)}\qquad &\psi(m(e(u),e(v))) = \fstar(u\concat v), \\
\text{(TB3)}\qquad &\psi(r(s)) = \psi(s)\qquad\text{for on-range }s\in\mathcal{S}.
\end{align*}

Approximate theorem-backedness replaces these exact equalities by explicit
budgets in expectation or on audited events.

\begin{proposition}[Broadest exact interface]
\label{prop:broadest-theorem-backed}
The exact theorem-backed interface is exactly the state-space analogue of the
three local laws. In particular, once \textup{(TB1)}--\textup{(TB3)} hold, all
preservation and optimizer-transfer results proved for OPS apply to the induced
reduction on $\mathcal{S}$ after replacing the textual oracle $\fstar$ by the
state semantics $\psi$.
\end{proposition}

\begin{proof}
The proof is the same induction as in Appendix~\ref{app:proofs-core}. The leaf
step uses \textup{(TB1)}. The parent step uses \textup{(TB2)}. Extra
re-summarization rounds use \textup{(TB3)}. No argument in the inductive
preservation proof depends on summaries being literal strings; it only uses the
fact that each realized leaf, merge, and on-range re-application preserves the
task value. Therefore the same chain of results goes through on $\mathcal{S}$.
\end{proof}

\paragraph{Two important special cases.}
\begin{itemize}[nosep]
    \item \textbf{Direct-summary route.} Take $\mathcal{S}=\Strings$,
    $e(x)=g(x)$, $m(s,t)=g(s\concat t)$, and $r(s)=g(s)$.
    \item \textbf{Codec route.} Take $\mathcal{S}$ to be a latent state space,
    let $e$ be an encoder, let $m$ be a learned or supplied state merge, and
    let $r$ be re-encode-after-decode on on-range states.
\end{itemize}

\subsection{Codec-Mediated Reductions}

The codec case is the mathematically clean way to unify learned sketches and
text-facing summarizers. The backend may carry an internal state, but the only
question is whether the supplied codec witnesses \textup{(TB1)}--\textup{(TB3)}.

\begin{proposition}[Supplied codec is enough]
\label{prop:codec-is-enough}
Suppose we are given a state space $\mathcal{S}$ with encoder
$e:\Strings\to\mathcal{S}$, decoder $d:\mathcal{S}\to\Strings$, state merge
$m:\mathcal{S}\times\mathcal{S}\to\mathcal{S}$, and semantic readout
$\psi:\mathcal{S}\to\Y$. Assume:
\begin{enumerate}[nosep]
    \item \emph{leaf preservation:} $\psi(e(x))=\fstar(x)$ for all $x$;
    \item \emph{merge compatibility:} $\psi(m(e(u),e(v)))=\fstar(u\concat v)$
    for all $u,v$;
    \item \emph{re-encoding stability:} $\psi(e(d(s)))=\psi(s)$ for every
    on-range state $s$.
\end{enumerate}
Then the induced reduction with leaf map $e$, merge map $m$, and re-summary map
$r(s):=e(d(s))$ is exact theorem-backed.
\end{proposition}

\begin{proof}
The three assumptions are exactly \textup{(TB1)}--\textup{(TB3)} after setting
$r(s)=e(d(s))$. So Proposition~\ref{prop:broadest-theorem-backed} applies
immediately. The point is that we do not need a theorem-domain autoencoder to
be \emph{learned inside the backend}. We only need a supplied codec with these
properties.
\end{proof}

\begin{remark}[Why this unifies LLM and neural routes]
The direct-summary route is the codec route with the identity decoder.
Conversely, a learned sketch or latent-state model becomes theorem-backed as
soon as someone supplies a state semantics and a codec witness satisfying
Proposition~\ref{prop:codec-is-enough}. So the distinction is not between
``LLM'' and ``neural'' systems. It is between operators that do and do not come
with a theorem-domain witness.
\end{remark}

\subsection{Relation to the Global $A1/A2/A3$ Route}

Appendix~\ref{sec:global} introduced the deterministic global assumptions:
global sufficiency (Assumption~\ref{ass:app1}), global two-route compatibility
(Assumption~\ref{ass:app2}), and oracle-level dependence of on-range merges
(Assumption~\ref{ass:app3}).

\begin{proposition}[Exact all-input collapse to $A1/A2/A3$]
\label{prop:all-input-collapse}
For deterministic direct-summary reductions, exact theorem-backedness on all
one-leaf and two-leaf inputs together with Assumption~\ref{ass:app3} is
equivalent to the full global regime
\textup{Assumptions~\ref{ass:app1}--\ref{ass:app3}}.
\end{proposition}

\begin{proof}
Assume first that Assumptions~\ref{ass:app1}--\ref{ass:app3} hold. Then global
sufficiency is exactly the leaf law for every input, and global two-route
compatibility is exactly the merge law for every pair of inputs. The
re-summary/on-range dependence part is Assumption~\ref{ass:app3}, so the exact
theorem-backed conditions hold everywhere.

Conversely, suppose the direct reduction is exact theorem-backed on every
one-leaf and two-leaf input, and assume Assumption~\ref{ass:app3}. Applying the
one-leaf statement to an arbitrary $z\in\Strings$ gives
Assumption~\ref{ass:app1}. Applying the two-leaf statement to an arbitrary pair
$(u,v)$ gives Assumption~\ref{ass:app2}. Together with the assumed
Assumption~\ref{ass:app3}, this is exactly the global $A1/A2/A3$ regime.
\end{proof}

So the local-law bundle is the broadest theorem-backed interface, while the
global algebraic regime is a stronger deterministic specialization.

\subsection{Mergeable Sketches as a Special Case}

Appendix~\ref{app:mergeable} explained the connection to classical mergeable
summaries. The theorem-backed perspective sharpens that link.

\begin{proposition}[Exact mergeable sketch plus on-range stability]
\label{prop:mergeable-sketch-special-case}
Let $\mathcal{S}$ be a deterministic sketch space with encoder
$e:\Strings\to\mathcal{S}$, merge $\oplus$, decoder $d$, and readout
$\psi:\mathcal{S}\to\Y$. Suppose:
\begin{enumerate}[nosep]
    \item $\psi(e(x))=\fstar(x)$ for all leaves $x$;
    \item $\psi(e(u)\oplus e(v))=\fstar(u\concat v)$ for all pairs $(u,v)$;
    \item $\psi(e(d(s)))=\psi(s)$ for every on-range sketch state $s$.
\end{enumerate}
Then the sketch is an exact theorem-backed reduction. If the task is
commutative and the global assumptions of Appendix~\ref{sec:global} hold, this
reduces to the classical mergeable-summary picture.
\end{proposition}

\begin{proof}
This is Proposition~\ref{prop:codec-is-enough} with $m=\oplus$. The extra
on-range stability condition is exactly the additional ingredient classical
mergeability does not provide by itself. It is the OPS-style idempotence
requirement needed for safe repeated summary calls. In the commutative exact
regime, Appendix~\ref{app:mergeable} shows that the resulting operator is
precisely a mergeable summary in the classical sense.
\end{proof}

\paragraph{Takeaway.}
The right top-level object is not ``an LLM summarizer'' or ``a neural sketch.''
It is a compositional reduction operator equipped with either exact local laws,
approximate local laws, or a stronger route such as
Assumptions~\ref{ass:app1}--\ref{ass:app3} that compiles into those laws.
